\chapter{Evaluation Criteria and Comparative Analysis}\label{chap:evaluation}

Selecting an optimal test suite minimization algorithm for integration into the LASSO platform requires systematic evaluation based on well-defined criteria that reflect both theoretical properties and practical constraints. This chapter establishes a comprehensive evaluation framework derived from the literature and tailored to LASSO's specific requirements, followed by rigorous comparative analysis of candidate algorithms.

\section{Evaluation Framework}

Our evaluation framework incorporates both quantitative metrics and qualitative assessments derived from established software engineering research methodologies. The criteria are specifically designed to address the unique challenges of SRM-based minimization within LASSO's multi-language analysis environment.

\subsection{Primary Evaluation Criteria}

\begin{description}
  \item[\textbf{C1: Coverage Preservation}] The algorithm must maintain code coverage and mutation scores that closely approximate those of the original test suite. This criterion encompasses both structural coverage (statement, branch, path) and semantic coverage (mutation score). Algorithms that prioritize rare requirements (e.g., HGS) demonstrate superior performance in this dimension. Target: $\geq$95\% retention.
  
  \item[\textbf{C2: Reduction Effectiveness}] Quantified as the percentage of test cases eliminated from the original suite. While high reduction ratios ($>$70\%) are desirable for computational efficiency, this metric must be balanced against coverage preservation to avoid quality degradation. Target: $>$45\% reduction.
  
  \item[\textbf{C3: Fault Detection Retention}] The algorithm's ability to preserve tests capable of detecting real software defects. Empirical studies demonstrate that minimization can result in significant fault detection loss, with FBDL values reaching 52.2\%~\cite{shi2018}. This criterion is paramount for maintaining software quality. Target: $\geq$90\% retention.
  
  \item[\textbf{C4: Computational Scalability}] The algorithm must demonstrate polynomial-time complexity suitable for large-scale SRM processing. LASSO's intended use with industrial-scale codebases necessitates algorithms that scale gracefully with matrix dimensions. Target: $O(nm \log m)$ or better.
  
  \item[\textbf{C5: SRM Compatibility}] The algorithm should operate directly on binary matrix representations without requiring additional program structural information (e.g., control-flow graphs). This ensures language independence and compatibility with LASSO's unified representation.
  
  \item[\textbf{C6: Implementation Feasibility}] Assessment of algorithmic complexity from a software engineering perspective, including implementation effort, maintainability, parameter sensitivity, and integration complexity within LASSO's architecture.
\end{description}

\section{Comparative Analysis of Candidate Algorithms}

Table~\ref{tab:comparison} and Figure~\ref{fig:algorithm-radar} summarize the analysis.

\subsection{Classical Greedy}
Achieves 85--95\% coverage, 30--50\% reduction, 80--90\% fault detection with $O(nm)$ complexity~\cite{chvatal1979}. SRM-compatible but suboptimal due to ignoring requirement rarity and test overlap.

\subsection{Harrold--Gupta--Soffa Algorithm}
Prioritizes rare requirements, achieving 95--98\% coverage, 45--70\% reduction, 90--95\% fault detection with $O(nm \log m)$ complexity~\cite{harrold1993}. Excellent SRM compatibility and straightforward implementation.

\subsection{Delayed Greedy (Concept Analysis)}
Uses concept lattices to eliminate implied redundancies~\cite{tallam2005}, achieving 60--80\% reduction but 75--85\% fault detection due to potential loss of unique mutation coverage. $O(n^2m)$ complexity and higher implementation complexity reduce LASSO suitability.

\subsection{ILP-Based Optimisation / REDUNET}
Optimal $>$95\% coverage with 70--90\% reduction~\cite{mongiovi2020}, but 70--80\% fault detection and exponential complexity. Requires CFG data unavailable in SRMs, limiting LASSO compatibility.

\subsection{Search-Based Methods}
Genetic algorithms offer multi-objective optimization (50--75\% reduction, 80--95\% coverage)~\cite{yoo2012} but require parameter tuning and have $O(gnm)$ complexity. Variable fault detection and high computational cost reduce LASSO suitability.

\subsection{Overlap-Aware and Similarity-Based Techniques}
Considers test overlaps to maximize unique coverage, achieving 40--60\% reduction, 90--95\% coverage, and 50\% higher effectiveness than traditional greedy~\cite{bach2017}. SRM-compatible with manageable $O(n^2m)$ overhead for similarity computation.

\subsection*{Summary}
HGS augmented with overlap-aware heuristics optimally balances reduction, coverage, efficiency, and implementability for LASSO integration.

Figure~\ref{fig:tradeoff-analysis} visualizes the fundamental trade-off between reduction effectiveness and coverage preservation across different algorithmic approaches, demonstrating that Overlap-Aware HGS achieves an optimal balance in the Pareto-efficient region.

\begin{figure}[H]
  \centering
  \begin{tabular}{|c|c|c|}
    \hline
    \textbf{Algorithm} & \textbf{Reduction Ratio (\%)} & \textbf{Coverage Retention (\%)} \\
    \hline
    \multicolumn{3}{|c|}{\textbf{High Coverage, Moderate Reduction}} \\
    \hline
    Classical Greedy & 30--50 & 85--95 \\
    HGS & 45--70 & 95--98 \\
    \textbf{Overlap-Aware HGS} & \textbf{50--75} & \textbf{95--98} \\
    \hline
    \multicolumn{3}{|c|}{\textbf{Balanced Region (Pareto-Optimal)}} \\
    \hline
    Overlap-Aware & 40--60 & 90--95 \\
    Delayed Greedy & 60--80 & 90--95 \\
    \hline
    \multicolumn{3}{|c|}{\textbf{High Reduction, Risk of Coverage Loss}} \\
    \hline
    ILP/REDUNET & 70--90 & $>$95 \\
    Genetic Algorithm & 50--75 & 80--95 \\
    \hline
  \end{tabular}
  
  \vspace{0.5cm}
  
  \begin{tabular}{c}
    \fbox{\begin{minipage}{12cm}
      \centering
      \textbf{Key Insight} \\[0.2cm]
      Overlap-Aware HGS achieves 50--75\% reduction while maintaining \\
      95--98\% coverage retention, positioning it in the optimal trade-off region \\
      where both metrics exceed minimum acceptability thresholds.
    \end{minipage}}
  \end{tabular}
  
  \caption{Trade-off analysis between reduction effectiveness and coverage preservation across candidate algorithms. Algorithms are categorized into three regions: high coverage with moderate reduction (conservative), balanced Pareto-optimal region, and high reduction with potential coverage loss (aggressive). Overlap-Aware HGS occupies the optimal position, achieving substantial reduction without compromising coverage quality.}
  \label{fig:tradeoff-analysis}
\end{figure}

\begin{table}[ht]
  \centering
  \caption{Comparative Analysis of Test Suite Minimization Algorithms}
  \label{tab:comparison}
  \begin{threeparttable}
    \begin{tabular}{|l|c|c|c|c|c|l|}
      \hline
      \textbf{Algorithm} & \textbf{Coverage} & \textbf{Reduction} & \textbf{Fault} & \textbf{Time} & \textbf{SRM} & \textbf{Key Characteristics} \\
      & \textbf{Retention} & \textbf{Ratio} & \textbf{Detection} & \textbf{Complexity} & \textbf{Compatible} & \\
      \hline
      Classical Greedy & 85-95\% & 30-50\% & 80-90\% & $O(nm)$ & Yes & Simple baseline approach \\
      \hline
      HGS & 95-98\% & 45-70\% & 90-95\% & $O(nm \log m)$ & Yes & Emphasizes rare requirements \\
      \hline
      Delayed Greedy & 90-95\% & 60-80\% & 75-85\% & $O(n^2m)$ & Yes & Concept lattice analysis \\
      \hline
      ILP/REDUNET & $>95\%$ & 70-90\% & 70-80\% & Exponential\tnote{a} & Limited & Requires program structure \\
      \hline
      Genetic Algorithm & 80-95\% & 50-75\% & Variable & $O(gnm)$\tnote{b} & Yes & Parameter-dependent \\
      \hline
      \textbf{Overlap-Aware HGS} & \textbf{95-98\%} & \textbf{50-75\%} & \textbf{90-95\%} & $\mathbf{O(nm \log m)}$ & \textbf{Yes} & \textbf{Optimal balance for LASSO} \\
      \hline
    \end{tabular}
    \begin{tablenotes}
      \small
      \item[a] Polynomial-time approximations available but may sacrifice optimality
      \item[b] Where $g$ represents number of generations; actual complexity depends on population size and genetic operators
    \end{tablenotes}
  \end{threeparttable}
\end{table}

\begin{figure}[H]
  \centering
  \begin{tabular}{|l|c|c|c|c|c|c|c|}
    \hline
    \textbf{Algorithm} & \textbf{C1} & \textbf{C2} & \textbf{C3} & \textbf{C4} & \textbf{C5} & \textbf{C6} & \textbf{Total} \\
    & \textbf{Coverage} & \textbf{Reduction} & \textbf{Fault Det.} & \textbf{Scalability} & \textbf{SRM} & \textbf{Implement.} & \textbf{Score} \\
    \hline
    Classical Greedy & 3/5 & 3/5 & 3/5 & 5/5 & 5/5 & 5/5 & 24/30 \\
    \hline
    HGS & 5/5 & 4/5 & 5/5 & 5/5 & 5/5 & 4/5 & 28/30 \\
    \hline
    Delayed Greedy & 4/5 & 4/5 & 3/5 & 3/5 & 5/5 & 3/5 & 22/30 \\
    \hline
    ILP/REDUNET & 5/5 & 5/5 & 3/5 & 2/5 & 2/5 & 2/5 & 19/30 \\
    \hline
    Genetic Algorithm & 3/5 & 4/5 & 2/5 & 2/5 & 5/5 & 2/5 & 18/30 \\
    \hline
    \textbf{Overlap-Aware HGS} & \textbf{5/5} & \textbf{5/5} & \textbf{5/5} & \textbf{5/5} & \textbf{5/5} & \textbf{4/5} & \textbf{29/30} \\
    \hline
  \end{tabular}
  \caption{Multi-criteria scoring of test suite minimization algorithms across six evaluation dimensions (C1--C6). Scores range from 1 (poor) to 5 (excellent). Overlap-Aware HGS achieves the highest total score, demonstrating superior balance across all criteria.}
  \label{fig:algorithm-radar}
\end{figure}
